\documentclass[aspectratio=169]{beamer}
\usetheme{Madrid}
\usefonttheme{structurebold}
\usefonttheme[onlymath]{serif}

% Màu chủ đạo của template (xanh BK). Đổi mã RGB nếu muốn màu khác.
\definecolor{BKTemplate}{rgb}{0.11765, 0.30588, 0.47451}
\usecolortheme[named=BKTemplate]{structure}
\setbeamercolor{title}{fg=BKTemplate,bg=}
\setbeamertemplate{caption}[numbered]

%------------------------------------------------------------
% Gói lệnh (packages). Biên dịch bằng pdfLaTeX.
% Giữ nguyên bộ này để có sẵn công cụ cho hầu hết báo cáo.
%------------------------------------------------------------
% Lọc riêng cảnh báo vô hại của csquotes: encoding T5 (tiếng Việt) không có
% "space factor codes" cho bộ theo dõi dấu câu — không ảnh hưởng bản in.
\usepackage{silence}
\WarningFilter{csquotes}{No space factor codes}
\usepackage{amsmath}        % ký hiệu toán học, \mathbb ...
\usepackage{amsfonts}       % font toán, ký hiệu bảng đen
\usepackage{bm}             % \bm để bôi đậm ký hiệu Hy Lạp
\usepackage[utf8]{inputenc} % nhập UTF-8 (pdfLaTeX)
\usepackage[T5,T1]{fontenc} % T5 cho tiếng Việt (vntex), T1 cho Latin
\usepackage{lmodern}
\usepackage{url}            % gõ URL
\usepackage{booktabs}       % bảng đẹp (\toprule, \midrule ...)
\usepackage{nicefrac}       % phân số gọn 1/2
\usepackage[protrusion=false]{microtype} % vi chỉnh khoảng chữ (tắt protrusion:
                            % font T5/tiếng Việt không có bảng protrusion → tránh cảnh báo)
\usepackage{mathtools}
\usepackage{multicol}
\usepackage{comment}
\usepackage{forest}
\usepackage[linesnumbered, ruled]{algorithm2e}
\usepackage{algorithmic}
\usepackage{xcolor}
\usepackage{colortbl}
\usepackage{graphicx}
\usepackage{diagbox}
\usepackage{tikz}
\usepackage{subcaption}
\usepackage{dsfont}
\usepackage{stmaryrd}
\usepackage{csquotes}      % biblatex khuyến nghị (dùng chung với babel)
\usepackage[style=authortitle,backend=biber,language=english]{biblatex}
\addbibresource{references.bib}
% biblatex/csquotes không có gói bản địa hóa cho tiếng Việt → dùng định nghĩa
% tiếng Anh cho phần tài liệu tham khảo (các mục trong references.bib vốn tiếng Anh).
\DeclareLanguageMapping{vietnamese}{english}
% Hỗ trợ tiếng Việt (dấu, chia âm tiết). Đổi thành english nếu viết tiếng Anh.
\usepackage[vietnamese]{babel}
\DeclareQuoteAlias{english}{vietnamese} % dùng kiểu trích dẫn tiếng Anh cho ngữ cảnh tiếng Việt

% Khai báo metadata PDF thủ công để hyperref không phải suy ra từ \title/\author
% (tránh cảnh báo "Token not allowed in a PDF string" do \\ và ký tự tiếng Việt).
\hypersetup{pdftitle={Tieu de bao cao}, pdfauthor={Ho va ten}}

\renewcommand{\footnotesize}{\tiny}
\newcommand{\mathbi}[1]{\boldsymbol{#1}}
\newtheorem{proposition}{Proposition}

%------------------------------------------------------------
% THÔNG TIN TRANG BÌA — sửa các placeholder bên dưới
%------------------------------------------------------------
\title[Tiêu đề ngắn] % (tùy chọn) hiện ở thanh chân slide
{\texorpdfstring{\vspace*{1cm}\\\textbf{}\\}{}\textbf{Tiêu đề báo cáo của bạn}}

\subtitle{} % (tùy chọn) phụ đề

\author[Họ tên] % (tùy chọn) hiện ở thanh chân slide
{Supervisor: Tên giảng viên hướng dẫn\texorpdfstring{\\}{, }
Student: Họ và tên - MSSV}

\institute[HCMUT] % (tùy chọn)
{
  Faculty of Computer Science and Engineering\\
  Ho Chi Minh City University of Technology\\
  Vietnam National University Ho Chi Minh City
}

\date{\today} % hoặc thay bằng "Tháng MM, YYYY"

%------------------------------------------------------------
% In mục lục (TOC) ở đầu mỗi section và tô sáng section hiện tại.
% Đây là tính năng khung — giữ nguyên.
%------------------------------------------------------------
\AtBeginSection[]
{
  \begin{frame}
    \frametitle{Table of Contents}
    \tableofcontents[currentsection]
  \end{frame}
}

% Lệnh tiện ích: khung ảnh placeholder vẽ bằng TikZ (dùng khi chưa có hình thật).
% Thay \figplaceholder{...}{...} bằng \includegraphics{images/ten-anh} khi có ảnh.
\newcommand{\figplaceholder}[2]{%
  \begin{tikzpicture}
    \draw[thick, dashed, gray, fill=gray!10] (0,0) rectangle (#1,#2);
    \node[gray] at (#1/2,#2/2) {\itshape Figure placeholder};
  \end{tikzpicture}%
}

\begin{document}

%------------------------------------------------------------
% TRANG BÌA (dùng ảnh nền FrontBackground.pdf)
%------------------------------------------------------------
\setbeamertemplate{background canvas}{
	\includegraphics[width=\paperwidth,height=\paperheight]{images/FrontBackground.pdf}
}
\begin{frame}[plain]
\titlepage
\end{frame}

% Đặt ảnh nền Background.pdf cho tất cả các slide còn lại.
\setbeamertemplate{background canvas}{
	\includegraphics[width=\paperwidth,height=\paperheight]{images/Background.pdf}
}

%------------------------------------------------------------
% Mục lục sau trang bìa
%------------------------------------------------------------
\begin{frame}
\frametitle{Table of Contents}
\tableofcontents
\end{frame}

%============================================================
% ===== SLIDE MẪU: sao chép / xóa / sửa tùy nhu cầu =====
%============================================================

%---------------------------------------------------------
\section{Section mẫu}
%---------------------------------------------------------

% -- SLIDE MẪU 1: danh sách (itemize) lồng nhau + alertblock --
\begin{frame}{Slide danh sách}
    \begin{itemize}
        \item \textbf{Ý chính thứ nhất}
        \begin{itemize}
            \item Ý phụ giải thích thêm.
            \item Ý phụ khác.
        \end{itemize}

        \vspace{0.3cm}

        \item \textbf{Ý chính thứ hai}
        \begin{itemize}
            \item Chi tiết bổ sung.
        \end{itemize}
    \end{itemize}

    \begin{alertblock}{Khối nhấn mạnh (alertblock)}
        Dùng để làm nổi bật một câu hỏi hoặc kết luận quan trọng.
    \end{alertblock}
\end{frame}

% -- SLIDE MẪU 2: khối (block) + công thức (equation) --
\begin{frame}{Slide công thức}
  \begin{block}{Khối nội dung (block)}
    \begin{itemize}
      \item Đặt định nghĩa, định lý hoặc diễn giải trong khối này.
      \item Ví dụ một phương trình:
      \begin{equation}
        f(x) = \sum_{i=1}^{n} w_i x_i + b
      \end{equation}
      trong đó $w_i$ là trọng số và $b$ là hệ số tự do.
    \end{itemize}
  \end{block}
\end{frame}

% -- SLIDE MẪU 3: hình ảnh (figure) --
% Khi có ảnh thật: thay \figplaceholder{8}{4} bằng
% \includegraphics[width=0.8\textwidth,keepaspectratio]{images/ten-anh.png}
\begin{frame}{Slide hình ảnh}
  \begin{figure}
    \centering
    \figplaceholder{8}{4}
    \caption{Chú thích cho hình minh họa.}
  \end{figure}
\end{frame}

% -- SLIDE MẪU 4: hai cột (columns) --
\begin{frame}{Slide hai cột}
  \begin{columns}
    \begin{column}{0.5\textwidth}
      \begin{itemize}
        \item Cột trái đặt văn bản / gạch đầu dòng.
        \item Phù hợp để mô tả kèm hình bên phải.
      \end{itemize}
    \end{column}
    \begin{column}{0.5\textwidth}
      \centering
      \figplaceholder{5}{3.5}
    \end{column}
  \end{columns}
\end{frame}

% -- SLIDE MẪU 5: bảng (booktabs) --
\begin{frame}{Slide bảng}
  \begin{table}
    \centering
    \begin{tabular}{lcc}
      \toprule
      \textbf{Phương pháp} & \textbf{Độ chính xác} & \textbf{Thời gian} \\
      \midrule
      Baseline   & 82.1\% & 1.0x \\
      Cải tiến A & 88.4\% & 1.2x \\
      Cải tiến B & \textbf{90.7\%} & 1.5x \\
      \bottomrule
    \end{tabular}
    \caption{Bảng so sánh mẫu (dùng gói booktabs).}
  \end{table}
\end{frame}

%============================================================
% ===== HẾT SLIDE MẪU =====
%============================================================

%---------------------------------------------------------
\section{References}
%---------------------------------------------------------
\begin{frame}
  \frametitle{Selected References}
  \nocite{*} % liệt kê mọi mục trong references.bib (bỏ dòng này nếu chỉ muốn mục đã \cite)
  \printbibliography
\end{frame}

%---------------------------------------------------------
% SLIDE KẾT THÚC
%---------------------------------------------------------
\setbeamertemplate{footline}{}
\begin{frame}
    \centering
    \vspace{1cm}
    \textbf{\Huge{- THE END - }}\\
    \vspace{0.2cm}
    \textit{\Large{Thank you for your attention}}\\
    \vspace{15pt}
    Acknowledgements:\\
    <lời cảm ơn của bạn>\\
    \vspace{20pt}
    Contact:\\
    <email>@hcmut.edu.vn
\end{frame}

\end{document}
